\chapter{Iterations}\label{ch:iterations}


Itteration 0: Florians model + his phyiscal contraption

What we wanted to do after looking at what already exists was adapt it for easier building by a user and since we have access to a bigger printer we wanted to avoid screws and prioritize ease of assembly.




itteration 1: has no camera support, size for tape holder is wrong, motor isnt mounted properly yet and we dont know what to do with then fan because we read something something pulling is better than pushing for consistency (source)


itteration 2: some sort of cover for the camera lense and fine tune the sizing for motor and tape roll, we still dont have a camera so camera size hole wasnt thought about at all just the closing lid for it


itteration 3: Finally got ahold of camera, so camera now slot into hole, made an oopsie with the lid because it doesnt account for the ribbon cable. motor gets attached slightly better with a gap so that gravity kinda holds it in place. and most importantly the dogshit lid mechanism we came up with is greatly simplified because friction holds it in place, no need to overengineer a solution when simple is better.


itteration 4: camera now slots in perfectly into the thingy, lid has a hole for ribbon cable, tape stays on, motor friction fits (but its sized wrongly so it makes the entire thing a bit bendy, version 5 will likely have that fixed. and the fan still doesnt have a mounting mechanism


itteration 5: TBD 
"""better motor mount, finally got a mosphet for the fan so some way to mount that would be good, led hole and connection"""


itteration 6: """we want proper microscopic test pictures like a hair or something."""


itteration 7: TBD
% """no room to mount the LED inside the box yet"""




\section{Web UI}

itteration 1: purely vibe coded ui for displaying camera and controlling the motor matches with the first three for the physical product


Itteration 2: same as before but with checks for if stuff is plugged in correctly, instructions for how to mount, and a mock camera that activates when no camera is connected. matches with 3-6 of the product itterations


itteration 3\label{sec:webui-3}

% Draft by Claude, edit freely.
The third version of the web interface was made for testing the camera (Section~\ref{sec:imaging-tests}).
To find good camera settings we had to take many pictures and keep track of how each one was taken.

\begin{itemize}
  \item \textbf{Faster live preview.} The preview used to update every 2~seconds. Shrinking the full
  12~megapixel picture took most of that time on the Raspberry Pi~3. The preview now skips pixels
  instead of resizing properly, and asks for the next picture as soon as the last one has arrived.
  One preview picture now takes about 0.3~s instead of about 1~s. A small progress bar shows when
  the next picture is on its way.
  \item \textbf{More camera settings.} Besides exposure and gain, the interface can now change white
  balance, brightness, contrast, saturation, sharpness, noise reduction, the time the LED is on before
  a picture, and whether to also save the raw sensor data. The settings are kept after a restart.
  \item \textbf{Notes on pictures.} Every picture can have a short note, for example about the light or
  the sample. The note, all settings, and what the sensor reported for that picture (actual exposure,
  gain, colour temperature, sensor temperature) are saved in a JSON file next to the picture. The file
  starts with a short explanation of its fields, so that other people, and AI tools, can read it
  without knowing the code.
  \item \textbf{Deleting pictures.} Getting a good picture takes many tries, so bad pictures can now be
  deleted together with their files.
\end{itemize}

% TODO: screenshot of the web interface (camera settings + recent pictures with notes)?
% Put it in figures/ and use \includegraphics[width=\linewidth]{webui-3}.








